% Formalization support lemmas.
%
% NOT book material. These declarations were invented in this workspace to
% carry a Lean formalization and have no counterpart in the source text --
% they carry no \dcref. They live here, outside the book chapters, so the
% book blueprint stays timeless mathematics while their \leanok marks and
% Lean anchors are preserved.
%
% Do not add book mathematics to this file, and do not add formalization
% notes to the book chapters.

\begin{exercise}
\label{prob:1-1}
\lean{problem1LineWithTwoOrigins_origins_nhds_inter}
\leanok
1-1. Let $X$ be the set of all points $\left( {x, y}\right)  \in  {\mathbb{R}}^{2}$ such that $y =  \pm  1$ , and let $M$ be the quotient of $X$ by the equivalence relation generated by $\left( {x, - 1}\right)  \sim  \left( {x,1}\right)$ for all $x \neq  0$ . Show that $M$ is locally Euclidean and second-countable, but not Hausdorff. (This space is called the line with two origins.)
\end{exercise}

\begin{exercise}
\label{prob:1-10}
\lean{standard_chart_matrix_unique}
\leanok
\uses{ex:1-36}
1-10. Let $k$ and $n$ be integers satisfying $0 < k < n$ , and let $P, Q \subseteq  {\mathbb{R}}^{n}$ be the linear subspaces spanned by $\left( {{e}_{1},\ldots ,{e}_{k}}\right)$ and $\left( {{e}_{k + 1},\ldots ,{e}_{n}}\right)$ , respectively, where ${e}_{i}$ is the $i$ th standard basis vector for ${\mathbb{R}}^{n}$ . For any $k$ -dimensional subspace $S \subseteq  {\mathbb{R}}^{n}$ that has trivial intersection with $Q$ , show that the coordinate representation $\varphi \left( S\right)$ constructed in Example 1.36 is the unique $\left( {n - k}\right)  \times  k$ matrix $B$ such that $S$ is spanned by the columns of the matrix $\left( \begin{matrix} {I}_{k} \\  B \end{matrix}\right)$ , where ${I}_{k}$ denotes the $k \times  k$ identity matrix.
\end{exercise}

\begin{exercise}
\label{prob:1-11}
\lean{closed_unit_ball_subtype_val_contMDiff}
\leanok
\uses{prob:1-7}
1-11. Let $M = {\overline{\mathbb{B}}}^{n}$ , the closed unit ball in ${\mathbb{R}}^{n}$ . Show that $M$ is a topological manifold with boundary in which each point in ${\mathbb{S}}^{n - 1}$ is a boundary point and each point in ${\mathbb{B}}^{n}$ is an interior point. Show how to give it a smooth structure such that every smooth interior chart is a smooth chart for the standard smooth structure on ${\mathbb{B}}^{n}$ . [Hint: consider the map $\pi  \circ  {\sigma }^{-1} : {\mathbb{R}}^{n} \rightarrow  {\mathbb{R}}^{n}$ , where $\sigma  : {\mathbb{S}}^{n} \rightarrow  {\mathbb{R}}^{n}$ is the stereographic projection (Problem 1-7) and $\pi$ is a projection from ${\mathbb{R}}^{n + 1}$ to ${\mathbb{R}}^{n}$ that omits some coordinate other than the last.]
\end{exercise}

\begin{exercise}
\label{prob:1-12}
\lean{prod_boundaryless_left_isManifold_with_boundary}
\leanok
\uses{prop:1-45}
1-12. Prove Proposition 1.45 (a product of smooth manifolds together with one smooth manifold with boundary is a smooth manifold with boundary).
\end{exercise}

\begin{exercise}
\label{prob:1-2}
\lean{Sigma.t2Space}
\mathlibok
1-2. Show that a disjoint union of uncountably many copies of $\mathbb{R}$ is locally Euclidean and Hausdorff, but not second-countable.
\end{exercise}

\begin{exercise}
\label{prob:1-4}
\lean{finite_option_bounding_family}
\leanok
1-4. Let $M$ be a topological manifold, and let $\mathcal{U}$ be an open cover of $M$ .

(a) Assuming that each set in $\mathcal{U}$ intersects only finitely many others, show that $\mathcal{U}$ is locally finite.

(b) Give an example to show that the converse to (a) may be false.

(c) Now assume that the sets in $\mathcal{U}$ are precompact in $M$ , and prove the converse: if $\mathcal{U}$ is locally finite, then each set in $\mathcal{U}$ intersects only finitely many others.
\end{exercise}

\begin{exercise}
\label{prob:1-5}
\lean{isSigmaCompact_connectedComponent_of_paracompact_t2_euclidean}
\leanok
1-5. Suppose $M$ is a locally Euclidean Hausdorff space. Show that $M$ is second-countable if and only if it is paracompact and has countably many connected components. [Hint: assuming $M$ is paracompact, show that each component of $M$ has a locally finite cover by precompact coordinate domains, and extract from this a countable subcover.]
\end{exercise}

\begin{exercise}
\label{prob:1-6}
\lean{center_chart_transition_eqOn_unit_ball}
\leanok
1-6. Let $M$ be a nonempty topological manifold of dimension $n \geq  1$ . If $M$ has a smooth structure, show that it has uncountably many distinct ones. [Hint: first show that for any $s > 0,{F}_{s}\left( x\right)  = {\left| x\right| }^{s - 1}x$ defines a homeomorphism from ${\mathbb{B}}^{n}$ to itself, which is a diffeomorphism if and only if $s = 1$ .]
\end{exercise}

\begin{exercise}
\label{prob:1-7}
\lean{stereographic_two_chart_same_smooth_structure}
\leanok
1-7. Let $N$ denote the north pole $\left( {0,\ldots ,0,1}\right)  \in  {\mathbb{S}}^{n} \subseteq  {\mathbb{R}}^{n + 1}$ , and let $S$ denote the south pole $\left( {0,\ldots ,0, - 1}\right)$ . Define the stereographic projection $\sigma  : {\mathbb{S}}^{n} \smallsetminus  \{ N\}  \rightarrow  {\mathbb{R}}^{n}$ by

\[
\sigma \left( {{x}^{1},\ldots ,{x}^{n + 1}}\right)  = \frac{\left( {x}^{1},\ldots ,{x}^{n}\right) }{1 - {x}^{n + 1}}.
\]

Let $\widetilde{\sigma }\left( x\right)  =  - \sigma \left( {-x}\right)$ for $x \in  {\mathbb{S}}^{n} \smallsetminus  \{ S\}$ .

(a) For any $x \in  {\mathbb{S}}^{n} \smallsetminus  \{ N\}$ , show that $\sigma \left( x\right)  = u$ , where $\left( {u,0}\right)$ is the point where the line through $N$ and $x$ intersects the linear subspace where ${x}^{n + 1} = 0$ (Fig. 1.13). Similarly, show that $\widetilde{\sigma }\left( x\right)$ is the point where the line through $S$ and $x$ intersects the same subspace. (For this reason, $\widetilde{\sigma }$ is called stereographic projection from the south pole.)

(b) Show that $\sigma$ is bijective, and

\[
{\sigma }^{-1}\left( {{u}^{1},\ldots ,{u}^{n}}\right)  = \frac{\left( 2{u}^{1},\ldots ,2{u}^{n},{\left| u\right| }^{2} - 1\right) }{{\left| u\right| }^{2} + 1}.
\]

(c) Compute the transition map $\widetilde{\sigma } \circ  {\sigma }^{-1}$ and verify that the atlas consisting of the two charts $\left( {{\mathbb{S}}^{n}\smallsetminus \{ N\} ,\sigma }\right)$ and $\left( {{\mathbb{S}}^{n}\smallsetminus \{ S\} ,\widetilde{\sigma }}\right)$ defines a smooth structure on ${\mathbb{S}}^{n}$ . (The coordinates defined by $\sigma$ or $\widetilde{\sigma }$ are called stereographic coordinates.)

(d) Show that this smooth structure is the same as the one defined in Example 1.31.

(Used on pp. 201, 269, 301, 345, 347, 450.)

![45_428_187_679_414_0.jpg](images/45_428_187_679_414_0.jpg)

Fig. 1.13 Stereographic projection
\end{exercise}

\begin{exercise}
\label{prob:1-9}
\lean{complexProjectiveSpaceIsManifold}
\leanok
1-9. Complex projective n-space, denoted by ${\mathbb{{CP}}}^{n}$ , is the set of all 1-dimensional complex-linear subspaces of ${\mathbb{C}}^{n + 1}$ , with the quotient topology inherited from the natural projection $\pi  : {\mathbb{C}}^{n + 1} \smallsetminus  \{ 0\}  \rightarrow  {\mathbb{{CP}}}^{n}$ . Show that ${\mathbb{{CP}}}^{n}$ is a compact ${2n}$ -dimensional topological manifold, and show how to give it a smooth structure analogous to the one we constructed for ${\mathbb{{RP}}}^{n}$ . (We use the correspondence

\[
\left( {{x}^{1} + i{y}^{1},\ldots ,{x}^{n + 1} + i{y}^{n + 1}}\right)  \leftrightarrow  \left( {{x}^{1},{y}^{1},\ldots ,{x}^{n + 1},{y}^{n + 1}}\right)
\]

to identify ${\mathbb{C}}^{n + 1}$ with ${\mathbb{R}}^{{2n} + 2}$ .) (Used on pp. 48,96,172,560,561.)
\end{exercise}

\begin{exercise}
\label{prob:2-1}
\lean{realHeavisideStep_not_smooth}
\leanok
2-1. Define $f : \mathbb{R} \rightarrow  \mathbb{R}$ by

\[
f\left( x\right)  = \left\{  \begin{array}{ll} 1, & x \geq  0 \\  0, & x < 0 \end{array}\right.
\]

Show that for every $x \in  \mathbb{R}$ , there are smooth coordinate charts $\left( {U,\varphi }\right)$ containing $x$ and $\left( {V,\psi }\right)$ containing $f\left( x\right)$ such that $\psi  \circ  f \circ  {\varphi }^{-1}$ is smooth as a map from $\varphi \left( {U \cap  {f}^{-1}\left( V\right) }\right)$ to $\psi \left( V\right)$ , but $f$ is not smooth in the sense we have defined in this chapter.
\end{exercise}

\begin{exercise}
\label{prob:2-13}
\lean{paracompactSpace_of_hasSubordinatePartitionOfUnity}
\leanok
2-13. Suppose $M$ is a topological space with the property that for every indexed open cover $\mathcal{X}$ of $M$ , there exists a partition of unity subordinate to $\mathcal{X}$ . Show that $M$ is paracompact.
\end{exercise}

\begin{exercise}
\label{prob:2-14}
\lean{exists_contMDiff_zero_iff_one_iff_of_isClosed}
\mathlibok
2-14. Suppose $A$ and $B$ are disjoint closed subsets of a smooth manifold $M$ . Show that there exists $f \in  {C}^{\infty }\left( M\right)$ such that $0 \leq  f\left( x\right)  \leq  1$ for all $x \in  M$ , ${f}^{-1}\left( 0\right)  = A$ , and ${f}^{-1}\left( 1\right)  = B$ .
\end{exercise}

\begin{exercise}
\label{prob:2-2}
\lean{contMDiff_pi_iff}
\leanok
\uses{prop:2-12}
2-2. Prove Proposition 2.12 (smoothness of maps into product manifolds).
\end{exercise}

\begin{exercise}
\label{prob:2-3}
\lean{hopf_map_smooth}
\leanok
2-3. For each of the following maps between spheres, compute sufficiently many coordinate representations to prove that it is smooth.

(a) ${p}_{n} : {\mathbb{S}}^{1} \rightarrow  {\mathbb{S}}^{1}$ is the nth power map for $n \in  \mathbb{Z}$ , given in complex notation by ${p}_{n}\left( z\right)  = {z}^{n}$ .

(b) $\alpha  : {\mathbb{S}}^{n} \rightarrow  {\mathbb{S}}^{n}$ is the antipodal map $\alpha \left( x\right)  =  - x$ .

(c) $F : {\mathbb{S}}^{3} \rightarrow  {\mathbb{S}}^{2}$ is given by $F\left( {w, z}\right)  = \left( {z\bar{w} + w\bar{z},{iw}\bar{z} - {iz}\bar{w}, z\bar{z} - w\bar{w}}\right)$ , where we think of ${\mathbb{S}}^{3}$ as the subset $\left\{  {\left( {w, z}\right)  : {\left| w\right| }^{2} + {\left| z\right| }^{2} = 1}\right\}$ of ${\mathbb{C}}^{2}$ .
\end{exercise}

\begin{exercise}
\label{prob:2-4}
\lean{closedUnitBall_exists_smoothManifoldWithBoundary}
\leanok
2-4. Show that the inclusion map ${\overline{\mathbb{B}}}^{n} \hookrightarrow  {\mathbb{R}}^{n}$ is smooth when ${\overline{\mathbb{B}}}^{n}$ is regarded as a smooth manifold with boundary.
\end{exercise}

\begin{exercise}
\label{prob:2-6}
\lean{contMDiff_homogeneousProjectivizationMap}
\leanok
2-6. Let $P : {\mathbb{R}}^{n + 1} \smallsetminus  \{ 0\}  \rightarrow  {\mathbb{R}}^{k + 1} \smallsetminus  \{ 0\}$ be a smooth function, and suppose that for some $d \in  \mathbb{Z}, P\left( {\lambda x}\right)  = {\lambda }^{d}P\left( x\right)$ for all $\lambda  \in  \mathbb{R} \smallsetminus  \{ 0\}$ and $x \in  {\mathbb{R}}^{n + 1} \smallsetminus  \{ 0\}$ . (Such a function is said to be homogeneous of degree $\mathbf{d}$ .) Show that the map $\widetilde{P} : {\mathbb{{RP}}}^{n} \rightarrow  {\mathbb{{RP}}}^{k}$ defined by $\widetilde{P}\left( \left\lbrack  x\right\rbrack  \right)  = \left\lbrack  {P\left( x\right) }\right\rbrack$ is well defined and smooth.
\end{exercise}

\begin{exercise}
\label{prob:2-7}
\lean{smooth_functions_not_finiteDimensional}
\leanok
2-7. Let $M$ be a nonempty smooth $n$ -manifold with or without boundary, and suppose $n \geq  1$ . Show that the vector space ${C}^{\infty }\left( M\right)$ is infinite-dimensional. [Hint: show that if ${f}_{1},\ldots ,{f}_{k}$ are elements of ${C}^{\infty }\left( M\right)$ with nonempty disjoint supports, then they are linearly independent.]
\end{exercise}

\begin{exercise}
\label{prob:2-8}
\lean{complexProjectiveAffineDiffeomorph}
\leanok
\uses{prob:1-9}
2-8. Define $F : {\mathbb{R}}^{n} \rightarrow  {\mathbb{{RP}}}^{n}$ by $F\left( {{x}^{1},\ldots ,{x}^{n}}\right)  = \left\lbrack  {{x}^{1},\ldots ,{x}^{n},1}\right\rbrack$ . Show that $F$ is a diffeomorphism onto a dense open subset of ${\mathbb{{RP}}}^{n}$ . Do the same for $G : {\mathbb{C}}^{n} \rightarrow  {\mathbb{{CP}}}^{n}$ defined by $G\left( {{z}^{1},\ldots ,{z}^{n}}\right)  = \left\lbrack  {{z}^{1},\ldots ,{z}^{n},1}\right\rbrack$ (see Problem 1-9).
\end{exercise}

\begin{exercise}
\label{prob:2-9}
\lean{exists_unique_smooth_complexProjectiveLine_polynomial_extension}
\leanok
\uses{prob:2-8}
2-9. Given a polynomial $p$ in one variable with complex coefficients, not identically zero, show that there is a unique smooth map $\widetilde{p} : {\mathbb{{CP}}}^{1} \rightarrow  {\mathbb{{CP}}}^{1}$ that makes the following diagram commute, where ${\mathbb{{CP}}}^{1}$ is 1-dimensional complex projective space and $G : \mathbb{C} \rightarrow  {\mathbb{{CP}}}^{1}$ is the map of Problem 2-8:

![63_708_297_231_215_0.jpg](images/63_708_297_231_215_0.jpg)

(Used on p. 465.)
\end{exercise}

\begin{exercise}
\label{prob:3-2}
\lean{tangentSpaceProdEquiv}
\leanok
\uses{prop:3-14}
3-2. Prove Proposition 3.14 (the tangent space to a product manifold).
\end{exercise}

\begin{exercise}
\label{prob:3-4}
\lean{tangentBundle_unitCircle_diffeomorph}
\leanok
3-4. Show that $T{\mathbb{S}}^{1}$ is diffeomorphic to ${\mathbb{S}}^{1} \times  \mathbb{R}$ .
\end{exercise}

\begin{exercise}
\label{prob:3-5}
\lean{not_exists_diffeomorph_maps_unitCircle_to_squareBoundary}
\leanok
3-5. Let ${\mathbb{S}}^{1} \subseteq  {\mathbb{R}}^{2}$ be the unit circle, and let $K \subseteq  {\mathbb{R}}^{2}$ be the boundary of the square of side 2 centered at the origin: $K = \left\{  {\left( {x, y}\right)  : \max \left( {\left| x\right| ,\left| y\right| }\right)  = 1}\right\}$ . Show that there is a homeomorphism $F : {\mathbb{R}}^{2} \rightarrow  {\mathbb{R}}^{2}$ such that $F\left( {\mathbb{S}}^{1}\right)  = K$ , but there is no diffeomorphism with the same property. [Hint: let $\gamma$ be a smooth curve whose image lies in ${\mathbb{S}}^{1}$ , and consider the action of ${dF}\left( {{\gamma }^{\prime }\left( t\right) }\right)$ on the coordinate functions $x$ and $y$ .] (Used on p. 123.)
\end{exercise}

\begin{exercise}
\label{prob:4-1}
\lean{euclideanHalfSpace_inclusion_not_diffeomorph_on_connected_open_neighborhoods_at_zero}
\leanok
\uses{thm:4-5}
4-1. Use the inclusion map ${\mathbb{H}}^{n} \hookrightarrow  {\mathbb{R}}^{n}$ to show that Theorem 4.5 does not extend to the case in which $M$ is a manifold with boundary. (Used on p. 80.)
\end{exercise}

\begin{exercise}
\label{prob:4-2}
\lean{contMDiffAt_map_mem_interior_of_nonsingular_mfderiv}
\leanok
4-2. Suppose $M$ is a smooth manifold (without boundary), $N$ is a smooth manifold with boundary, and $F : M \rightarrow  N$ is smooth. Show that if $p \in  M$ is a point such that $d{F}_{p}$ is nonsingular, then $F\left( p\right)  \in$ Int $N$ . (Used on pp. 80,87.)
\end{exercise}

\begin{exercise}
\label{prob:5-14}
\lean{existsUnique_diffeomorph_refl_of_isImmersedSubmanifold}
\leanok
\uses{thm:5-32}
5-14. Prove Theorem 5.32 (uniqueness of the smooth structure on an immersed submanifold once the topology is given).
\end{exercise}

\begin{exercise}
\label{prob:5-19}
\lean{figureEightCurveMap_counterexample}
\leanok
5-19. Suppose $S \subseteq  M$ is an embedded submanifold and $\gamma  : J \rightarrow  M$ is a smooth curve whose image happens to lie in $S$ . Show that ${\gamma }^{\prime }\left( t\right)$ is in the subspace ${T}_{\gamma \left( t\right) }S$ of ${T}_{\gamma \left( t\right) }M$ for all $t \in  J$ . Give a counterexample if $S$ is not embedded.
\end{exercise}

\begin{exercise}
\label{prob:5-20}
\lean{figureEightCurveMap_branch_velocity_not_tangent_to_other_branch}
\leanok
5-20. Show by giving a counterexample that the conclusion of Proposition 5.37 may be false if $S$ is merely immersed.
\end{exercise}

\begin{exercise}
\label{prob:5-4}
\lean{figureEightCurve_image_not_embedded_curve}
\leanok
\uses{ex:4-19}
5-4. Show that the image of the curve $\beta  : \left( {-\pi ,\pi }\right)  \rightarrow  {\mathbb{R}}^{2}$ of Example 4.19 is not an embedded submanifold of ${\mathbb{R}}^{2}$ . [Be careful: this is not the same as showing that $\beta$ is not an embedding.]
\end{exercise}

\begin{exercise}
\label{prob:6-3}
\lean{Manifold.exists_smooth_zero_on_and_positive_off_lt_of_isClosed}
\leanok
\uses{cor:6-22}
6-3. Let $M$ be a smooth manifold, let $B \subseteq  M$ be a closed subset, and let $\delta  : M \rightarrow \; \mathbb{R}$ be a positive continuous function. Show that there is a smooth function $\widetilde{\delta } : M \rightarrow  \mathbb{R}$ that is zero on $B$ , positive on $M \smallsetminus  B$ , and satisfies $\widetilde{\delta }\left( x\right)  < \delta \left( x\right)$ everywhere. [Hint: consider $f/\left( {f + 1}\right)$ , where $f$ is a smooth nonnegative function that vanishes exactly on $B$ , and use Corollary 6.22.]
\end{exercise}

\begin{exercise}
\label{prob:7-11}
\lean{hopf_action_isFree}
\leanok
7-11. Considering ${\mathbb{S}}^{{2n} + 1}$ as the unit sphere in ${\mathbb{C}}^{n + 1}$ , define an action of ${\mathbb{S}}^{1}$ on ${\mathbb{S}}^{{2n} + 1}$ , called the Hopf action, by

\[
z \cdot  \left( {{w}^{1},\ldots ,{w}^{n + 1}}\right)  = \left( {z{w}^{1},\ldots , z{w}^{n + 1}}\right) .
\]

Show that this action is smooth and its orbits are disjoint unit circles in ${\mathbb{C}}^{n + 1}$ whose union is ${\mathbb{S}}^{{2n} + 1}$ . (Used on p. 560.)
\end{exercise}

\begin{exercise}
\label{prob:7-12}
\lean{ContMDiffMonoidMorphism.toMulActionHom}
\leanok
7-12. Use the equivariant rank theorem to give another proof of Theorem 7.5 by showing that every Lie group homomorphism $F : G \rightarrow  H$ is equivariant with respect to suitable smooth $G$ -actions on $G$ and $H$ .
\end{exercise}

\begin{exercise}
\label{prob:7-17}
\lean{problem_7_17}
\leanok
7-17. Determine which of the following Lie groups are compact:

\[
\mathrm{{GL}}\left( {n,\mathbb{R}}\right) ,\mathrm{{SL}}\left( {n,\mathbb{R}}\right) ,\mathrm{{GL}}\left( {n,\mathbb{C}}\right) ,\mathrm{{SL}}\left( {n,\mathbb{C}}\right) ,\mathrm{U}\left( n\right) ,\mathrm{{SU}}\left( n\right) \text{ . }
\]
\end{exercise}

\begin{exercise}
\label{prob:7-2}
\lean{mfderiv_inv_at_one_apply}
\leanok
\uses{prop:3-14}
7-2. Let $G$ be a Lie group.

(a) Let $m : G \times  G \rightarrow  G$ denote the multiplication map. Using Proposition 3.14 to identify ${T}_{\left( e, e\right) }\left( {G \times  G}\right)$ with ${T}_{e}G \oplus  {T}_{e}G$ , show that the differential $d{m}_{\left( e, e\right) } : {T}_{e}G \oplus  {T}_{e}G \rightarrow  {T}_{e}G$ is given by

\[
d{m}_{\left( e, e\right) }\left( {X, Y}\right)  = X + Y
\]

[Hint: compute $d{m}_{\left( e, e\right) }\left( {X,0}\right)$ and $d{m}_{\left( e, e\right) }\left( {0, Y}\right)$ separately.]

(b) Let $i : G \rightarrow  G$ denote the inversion map. Show that $d{i}_{e} : {T}_{e}G \rightarrow  {T}_{e}G$ is given by $d{i}_{e}\left( X\right)  =  - X$ .

(Used on pp. 203, 522.)
\end{exercise}

\begin{exercise}
\label{prob:7-22}
\lean{quaternion_inv_eq_norm_sq_smul_star}
\leanok
7-22. Let $\mathbb{H} = \mathbb{C} \times  \mathbb{C}$ (considered as a real vector space), and define a bilinear product $\mathbb{H} \times  \mathbb{H} \rightarrow  \mathbb{H}$ by

\[
\left( {a, b}\right) \left( {c, d}\right)  = \left( {{ac} - d\bar{b},\bar{a}d + {cb}}\right) ,\;\text{ for }a, b, c, d \in  \mathbb{C}.
\]

With this product, $\mathbb{H}$ is a 4-dimensional algebra over $\mathbb{R}$ , called the algebra of quaternions. For each $p = \left( {a, b}\right)  \in  \mathbb{H}$ , define ${p}^{ * } = \left( {\bar{a}, - b}\right)$ . It is useful to work with the basis $\left( {\mathbb{1},\mathbb{i},\mathbb{j},\mathbb{k}}\right)$ for $\mathbb{H}$ defined by

\[
\mathbb{1} = \left( {1,0}\right) ,\;\mathbb{i} = \left( {i,0}\right) ,\;\mathbb{j} = \left( {0,1}\right) ,\;\mathbb{k} = \left( {0, - i}\right) .
\]

It is straightforward to verify that this basis satisfies

\[
{\widehat{\mathbb{i}}}^{2} = {\widehat{\mathbb{j}}}^{2} = {\mathbb{k}}^{2} =  - \mathbb{1},\;\mathbb{1}q = q\mathbb{1} = q\;\text{ for all }q \in  \mathbb{H},
\]

\[
\dot{\mathbb{i}}\widehat{j} =  - \widehat{\mathbb{j}}\widehat{i} = \mathbb{k},\;\dot{\mathbb{j}}\mathbb{k} =  - \mathbb{k}\widehat{j} = \widehat{\mathbb{i}},\;\mathbb{k}\widehat{\mathbb{i}} =  - \widehat{\mathbb{i}}\mathbb{k} = \widehat{\mathbb{j}},
\]

\[
{\mathbb{1}}^{ * } = \mathbb{1},\;{\dot{\mathbb{i}}}^{ * } =  - \dot{\mathbb{i}},\;{\dot{\mathbb{j}}}^{ * } =  - \dot{\mathbb{j}},\;{\mathbb{k}}^{ * } =  - \mathbb{k}.
\]

A quaternion $p$ is said to be real if ${p}^{ * } = p$ , and imaginary if ${p}^{ * } =  - p$ . Real quaternions can be identified with real numbers via the correspondence $x \leftrightarrow  x\mathbb{1}.$

(a) Show that quaternionic multiplication is associative but not commutative.

(b) Show that ${\left( pq\right) }^{ * } = {q}^{ * }{p}^{ * }$ for all $p, q \in  \mathbb{H}$ .

(c) Show that $\langle p, q\rangle  = \frac{1}{2}\left( {{p}^{ * }q + {q}^{ * }p}\right)$ is an inner product on $\mathbb{H}$ , whose associated norm satisfies $\left| {pq}\right|  = \left| p\right| \left| q\right|$ .

(d) Show that every nonzero quaternion has a two-sided multiplicative inverse given by ${p}^{-1} = {\left| p\right| }^{-2}{p}^{ * }$ .

(e) Show that the set ${\mathbb{H}}^{ * }$ of nonzero quaternions is a Lie group under quaternionic multiplication.

(Used on pp. 200, 200, 562.)
\end{exercise}

\begin{exercise}
\label{prob:7-23}
\lean{quaternionSphereContinuousMulEquivSpecialUnitary}
\leanok
\uses{prob:7-22}
7-23. Let ${\mathbb{H}}^{ * }$ be the Lie group of nonzero quaternions (Problem 7-22), and let $\mathcal{S} \subseteq  {\mathbb{H}}^{ * }$ be the set of unit quaternions. Show that $\mathcal{S}$ is a properly embedded Lie subgroup of ${\mathbb{H}}^{ * }$ , isomorphic to $\mathrm{{SU}}\left( 2\right)$ . (Used on pp. 200,562.)
\end{exercise}

\begin{exercise}
\label{prob:7-24}
\lean{tau_d_n_representation_faithful}
\leanok
7-24. Prove that each of the maps ${\tau }_{d}^{n} : \mathrm{{GL}}\left( {n,\mathbb{R}}\right)  \rightarrow  \mathrm{{GL}}\left( {\mathcal{P}}_{d}^{n}\right)$ described in Example ${7.36}\left( \mathrm{\;g}\right)$ is a faithful representation of $\mathrm{{GL}}\left( {n,\mathbb{R}}\right)$ .
\end{exercise}

\begin{exercise}
\label{prob:7-3}
\lean{lieGroup_of_contMDiff_mul}
\leanok
7-3. Our definition of Lie groups includes the requirement that both the multiplication map and the inversion map are smooth. Show that smoothness of the inversion map is redundant: if $G$ is a smooth manifold with a group structure such that the multiplication map $m : G \times  G \rightarrow  G$ is smooth, then $G$ is a Lie group. [Hint: show that the map $F : G \times  G \rightarrow  G \times  G$ defined by $F\left( {g, h}\right)  = \left( {g,{gh}}\right)$ is a bijective local diffeomorphism.]
\end{exercise}

\begin{exercise}
\label{prob:7-4}
\lean{generalLinear_det_fderiv_apply}
\leanok
\uses{cor:3-25}
7-4. Let det: $\operatorname{GL}\left( {n,\mathbb{R}}\right)  \rightarrow  \mathbb{R}$ denote the determinant function. Use Corollary 3.25 to compute the differential of det, as follows.

(a) For any $A \in  \mathbf{M}\left( {n,\mathbb{R}}\right)$ , show that

\[
{\left. \frac{d}{dt}\right| }_{t = 0}\det \left( {{I}_{n} + {tA}}\right)  = \operatorname{tr}A
\]

where $\operatorname{tr}\left( {A}_{j}^{i}\right)  = \mathop{\sum }\limits_{i}{A}_{i}^{i}$ is the trace of $A$ . [Hint: the defining equation (B.3) expresses $\det \left( {{I}_{n} + {tA}}\right)$ as a polynomial in $t$ . What is the linear term?]

(b) For $X \in  \mathrm{{GL}}\left( {n,\mathbb{R}}\right)$ and $B \in  {T}_{X}\mathrm{{GL}}\left( {n,\mathbb{R}}\right)  \cong  \mathrm{M}\left( {n,\mathbb{R}}\right)$ , show that

\[
d{\left( \det \right) }_{X}\left( B\right)  = \left( {\det X}\right) \operatorname{tr}\left( {{X}^{-1}B}\right) . \tag{7.11}
\]

[Hint: $\det \left( {X + {tB}}\right)  = \det \left( X\right) \det \left( {{I}_{n} + t{X}^{-1}B}\right)$ .]

(Used on p. 203.)
\end{exercise}

\begin{exercise}
\label{prob:7-5}
\lean{exists_lieGroupIsomorphism_of_universal_covering_group}
\leanok
\uses{thm:7-9}
7-5. Prove Theorem 7.9 (uniqueness of the universal covering group).
\end{exercise}

\begin{exercise}
\label{prob:7-6}
\lean{exists_nhds_one_subset_mul_inv_mem}
\leanok
7-6. Suppose $G$ is a Lie group and $U$ is any neighborhood of the identity. Show that there exists a neighborhood $V$ of the identity such that $V \subseteq  U$ and $g{h}^{-1} \in  U$ whenever $g, h \in  V$ . (Used on pp. 159,556.)
\end{exercise}

\begin{exercise}
\label{prob:7-8}
\lean{continuous_action_on_discrete_is_trivial}
\leanok
7-8. Suppose a connected topological group $G$ acts continuously on a discrete space $K$ . Show that the action is trivial. (Used on p. 562.)
\end{exercise}

\begin{exercise}
\label{prob:7-9}
\lean{real_projective_generalLinear_isPretransitive}
\leanok
7-9. Show that the formula

\[
A \cdot  \left\lbrack  x\right\rbrack   = \left\lbrack  {Ax}\right\rbrack
\]

defines a smooth, transitive left action of $\mathrm{{GL}}\left( {n + 1,\mathbb{R}}\right)$ on ${\mathbb{{RP}}}^{n}$ .
\end{exercise}

\begin{exercise}
\label{prob:8-1}
\lean{exists_supported_contMDiff_vectorField_extension_of_isClosed}
\leanok
\uses{lem:8-6}
8-1. Prove Lemma 8.6 (the extension lemma for vector fields).
\end{exercise}

\begin{exercise}
\label{prob:8-11}
\lean{problem_8_11_radius_squared_x_field_polar_coordinates}
\leanok
8-11. For each of the following vector fields on the plane, compute its coordinate representation in polar coordinates on the right half-plane $\{ \left( {x, y}\right)  : x > 0\}$ .

(a) $X = x\frac{\partial }{\partial x} + y\frac{\partial }{\partial y}$ .

(b) $Y = x\frac{\partial }{\partial y} - y\frac{\partial }{\partial x}$ .

(c) $Z = \left( {{x}^{2} + {y}^{2}}\right) \frac{\partial }{\partial x}$ .
\end{exercise}

\begin{exercise}
\label{prob:8-16}
\lean{problem_8_16_lie_bracket}
\leanok
8-16. For each of the following pairs of vector fields $X, Y$ defined on ${\mathbb{R}}^{3}$ , compute the Lie bracket $\left\lbrack  {X, Y}\right\rbrack$ .

(a) $X = y\frac{\partial }{\partial z} - {2x}{y}^{2}\frac{\partial }{\partial y};\;Y = \frac{\partial }{\partial y}$ .

(b) $X = x\frac{\partial }{\partial y} - y\frac{\partial }{\partial x};\;Y = y\frac{\partial }{\partial z} - z\frac{\partial }{\partial y}$ .

(c) $X = x\frac{\partial }{\partial y} - y\frac{\partial }{\partial x};\;Y = x\frac{\partial }{\partial y} + y\frac{\partial }{\partial x}.$
\end{exercise}

\begin{exercise}
\label{prob:8-19}
\lean{Cross.lieAlgebra}
\leanok
8-19. Show that ${\mathbb{R}}^{3}$ with the cross product is a Lie algebra.
\end{exercise}

\begin{exercise}
\label{prob:8-2}
\lean{euler_homogeneous_function_theorem}
\leanok
\uses{ex:8-3}
8-2. EULER's HOMOGENEOUS FUNCTION THEOREM: Let $c$ be a real number, and let $f : {\mathbb{R}}^{n} \smallsetminus  \{ 0\}  \rightarrow  \mathbb{R}$ be a smooth function that is positively homogeneous of degree $c$ , meaning that $f\left( {\lambda x}\right)  = {\lambda }^{c}f\left( x\right)$ for all $\lambda  > 0$ and $x \in  {\mathbb{R}}^{n} \smallsetminus  \{ 0\}$ . Prove that ${Vf} = {cf}$ , where $V$ is the Euler vector field defined in Example 8.3. (Used on p. 248.)
\end{exercise}

\begin{exercise}
\label{prob:8-20}
\lean{problem_8_20_cross_equiv_A}
\leanok
8-20. Let $A \subseteq  \mathfrak{X}\left( {\mathbb{R}}^{3}\right)$ be the subspace spanned by $\{ X, Y, Z\}$ , where

\[
X = y\frac{\partial }{\partial z} - z\frac{\partial }{\partial y},\;Y = z\frac{\partial }{\partial x} - x\frac{\partial }{\partial z},\;Z = x\frac{\partial }{\partial y} - y\frac{\partial }{\partial x}.
\]

Show that $A$ is a Lie subalgebra of $\mathfrak{X}\left( {\mathbb{R}}^{3}\right)$ , which is isomorphic to ${\mathbb{R}}^{3}$ with the cross product. (Used on p. 538.)
\end{exercise}

\begin{exercise}
\label{prob:8-21}
\lean{gl2_real_two_dimensional_models_not_lie_equiv}
\leanok
8-21. Prove that up to isomorphism, there are exactly one 1-dimensional Lie algebra and two 2-dimensional Lie algebras. Show that all three algebras are isomorphic to Lie subalgebras of $\mathfrak{{gl}}\left( {2,\mathbb{R}}\right)$ .
\end{exercise}

\begin{exercise}
\label{prob:8-22}
\lean{algebra_derivations_lie_subalgebra}
\leanok
8-22. Let $A$ be any algebra over $\mathbb{R}$ . A derivation of $A$ is a linear map $D : A \rightarrow  A$ satisfying $D\left( {xy}\right)  = \left( {Dx}\right) y + x\left( {Dy}\right)$ for all $x, y \in  A$ . Show that if ${D}_{1}$ and ${D}_{2}$ are derivations of $A$ , then $\left\lbrack  {{D}_{1},{D}_{2}}\right\rbrack   = {D}_{1} \circ  {D}_{2} - {D}_{2} \circ  {D}_{1}$ is also a derivation. Show that the set of derivations of $A$ is a Lie algebra with this bracket operation.
\end{exercise}

\begin{exercise}
\label{prob:8-23}
\lean{GroupLieAlgebra.prodLieEquiv}
\leanok
8-23. (a) Given Lie algebras $\mathfrak{g}$ and $\mathfrak{h}$ , show that the direct sum $\mathfrak{g} \oplus  \mathfrak{h}$ is a Lie algebra with the bracket defined by

\[
\left\lbrack  {\left( {X, Y}\right) ,\left( {{X}^{\prime },{Y}^{\prime }}\right) }\right\rbrack   = \left( {\left\lbrack  {X,{X}^{\prime }}\right\rbrack  ,\left\lbrack  {Y,{Y}^{\prime }}\right\rbrack  }\right) .
\]

(b) Suppose $G$ and $H$ are Lie groups. Prove that $\operatorname{Lie}\left( {G \times  H}\right)$ is isomorphic to $\operatorname{Lie}\left( G\right)  \oplus  \operatorname{Lie}\left( H\right)$ .
\end{exercise}

\begin{exercise}
\label{prob:8-24}
\lean{inversion_pushforward_on_groupLieAlgebra_map_lie}
\leanok
8-24. Suppose $G$ is a Lie group and $\mathfrak{g}$ is its Lie algebra. A vector field $X \in  \mathfrak{X}\left( G\right)$ is said to be right-invariant if it is invariant under all right translations.

(a) Show that the set $\overline{\mathrm{g}}$ of right-invariant vector fields on $G$ is a Lie subalgebra of $\mathfrak{X}\left( G\right)$ .

(b) Let $i : G \rightarrow  G$ denote the inversion map $i\left( g\right)  = {g}^{-1}$ . Show that the push-forward ${i}_{ * } : \mathcal{X}\left( G\right)  \rightarrow  \mathcal{X}\left( G\right)$ restricts to a Lie algebra isomorphism from $\mathfrak{g}$ to $\overline{\mathfrak{g}}$ .
\end{exercise}

\begin{exercise}
\label{prob:8-25}
\lean{AddGroupLieAlgebra.instIsLieAbelian_of_addCommGroup}
\leanok
\uses{prob:7-2}
8-25. Prove that if $G$ is an abelian Lie group, then $\operatorname{Lie}\left( G\right)$ is abelian. [Hint: show that the inversion map $i : G \rightarrow  G$ is a group homomorphism, and use Problem 7-2.]
\end{exercise}

\begin{exercise}
\label{prob:8-3}
\lean{smoothVectorField_not_finiteDimensional}
\leanok
8-3. Let $M$ be a nonempty positive-dimensional smooth manifold with or without boundary. Show that $\mathfrak{X}\left( M\right)$ is infinite-dimensional.
\end{exercise}

\begin{exercise}
\label{prob:8-31}
\lean{LieHom.mem_ker}
\mathlibok
8-31. Let $\mathfrak{g}$ be a Lie algebra. A linear subspace $\mathfrak{h} \subseteq  \mathfrak{g}$ is called an ideal in $\mathfrak{g}$ if $\left\lbrack  {X, Y}\right\rbrack   \in  \mathfrak{h}$ whenever $X \in  \mathfrak{h}$ and $Y \in  \mathfrak{g}.$

(a) Show that if $\mathfrak{h}$ is an ideal in $\mathfrak{g}$ , then the quotient space $\mathfrak{g}/\mathfrak{h}$ has a unique Lie algebra structure such that the projection $\pi  : \mathfrak{g} \rightarrow  \mathfrak{g}/\mathfrak{h}$ is a Lie algebra homomorphism.

(b) Show that a subspace $\mathfrak{h} \subseteq  \mathfrak{g}$ is an ideal if and only if it is the kernel of a Lie algebra homomorphism.

(Used on p. 533.)
\end{exercise}

\begin{exercise}
\label{prob:8-5}
\lean{exists_localFrame_extension_of_closed}
\leanok
\uses{prop:8-11}
8-5. Prove Proposition 8.11 (completion of local frames).
\end{exercise}

\begin{exercise}
\label{prob:8-6}
\lean{problem_8_6_X3_coordinates}
\leanok
\uses{prob:7-22,prob:7-23}
8-6. Let $\mathbb{H}$ be the algebra of quaternions and let $\mathcal{S} \subseteq  \mathbb{H}$ be the group of unit quaternions (see Problems 7-22 and 7-23).

(a) Show that if $p \in  \mathbb{H}$ is imaginary, then ${qp}$ is tangent to $\mathcal{S}$ at each $q \in  \mathcal{S}$ . (Here we are identifying each tangent space to $\mathbb{H}$ with $\mathbb{H}$ itself in the usual way.)

(b) Define vector fields ${X}_{1},{X}_{2},{X}_{3}$ on $\mathbb{H}$ by

\[
{\left. {X}_{1}\right| }_{q} = q{\left. {\left. {\left. {X}_{2}\right| }_{q} = q\right| }_{p},\;{X}_{3}\right| }_{q} = q\mathbb{k}.
\]

Show that these vector fields restrict to a smooth left-invariant global frame on $S$ .

(c) Under the isomorphism $\left( {{x}^{1},{x}^{2},{x}^{3},{x}^{4}}\right)  \leftrightarrow  {x}^{1}\mathbb{1} + {x}^{2}\mathbb{i} + {x}^{3}\mathbb{j} + {x}^{4}\mathbb{k}$ between ${\mathbb{R}}^{4}$ and $\mathbb{H}$ , show that these vector fields have the following coordinate representations:

\[
{X}_{1} =  - {x}^{2}\frac{\partial }{\partial {x}^{1}} + {x}^{1}\frac{\partial }{\partial {x}^{2}} + {x}^{4}\frac{\partial }{\partial {x}^{3}} - {x}^{3}\frac{\partial }{\partial {x}^{4}},
\]

\[
{X}_{2} =  - {x}^{3}\frac{\partial }{\partial {x}^{1}} - {x}^{4}\frac{\partial }{\partial {x}^{2}} + {x}^{1}\frac{\partial }{\partial {x}^{3}} + {x}^{2}\frac{\partial }{\partial {x}^{4}},
\]

\[
{X}_{3} =  - {x}^{4}\frac{\partial }{\partial {x}^{1}} + {x}^{3}\frac{\partial }{\partial {x}^{2}} - {x}^{2}\frac{\partial }{\partial {x}^{3}} + {x}^{1}\frac{\partial }{\partial {x}^{4}}.
\]

(Used on pp. 179, 562.)
\end{exercise}

\begin{exercise}
\label{prob:8-7}
\lean{problem_8_7_isLocalFrameOn}
\leanok
\uses{prob:7-22,prob:8-6}
8-7. The algebra of octonions (also called Cayley numbers) is the 8-dimensional real vector space $\mathbb{O} = \mathbb{H} \times  \mathbb{H}$ (where $\mathbb{H}$ is the space of quaternions defined in Problem 7-22) with the following bilinear product:

\[
\left( {p, q}\right) \left( {r, s}\right)  = \left( {{pr} - s{q}^{ * },{p}^{ * }s + {rq}}\right) ,\;\text{ for }p, q, r, s \in  \mathbb{H}. \tag{8.20}
\]

Show that $\mathbb{O}$ is a noncommutative, nonassociative algebra over $\mathbb{R}$ , and prove that there exists a smooth global frame on ${\mathbb{S}}^{7}$ by imitating as much of Problem 8-6 as you can. [Hint: it might be helpful to prove that $\left( {P{Q}^{ * }}\right) Q = \; P\left( {{Q}^{ * }Q}\right)$ for all $P, Q \in  \mathbb{O}$ , where ${\left( p, q\right) }^{ * } = \left( {{p}^{ * }, - q}\right)$ . For more about the octonions, see [Bae02].] (Used on p. 179.)
\end{exercise}

\begin{exercise}
\label{prob:8-8}
\lean{problem_8_8_not_linearIndependent_of_mfderiv_eq}
\leanok
8-8. The algebra of sedenions is the 16-dimensional real vector space $\mathbb{S} = \mathbb{O} \times  \mathbb{O}$ with the product defined by (8.20), but with $p, q, r$ , and $s$ interpreted as elements of $\mathbb{O}$ . Why does sedenionic multiplication not yield a global frame for ${\mathbb{S}}^{15}$ ? [Remark: the name "sedenion" comes from the Latin sedecim, meaning sixteen. A division algebra is an algebra with a multiplicative identity element and no zero divisors (i.e., ${ab} = 0$ if and only if $a = 0$ or $b = 0$ ). It follows from the work of Bott, Milnor, and Kervaire on parallelizability of spheres [MB58, Ker58] that a finite-dimensional division algebra over $\mathbb{R}$ must have dimension 1, 2, 4, or 8.]
\end{exercise}

\begin{exercise}
\label{prob:8-9}
\lean{problem_8_9_no_smooth_related_vector_field}
\leanok
\uses{prop:8-19}
8-9. Show by finding a counterexample that Proposition 8.19 is false if we replace the assumption that $F$ is a diffeomorphism by the weaker assumption that it is smooth and bijective.
\end{exercise}
